                  IMC 2014, Blagoevgrad, Bulgaria


                              Day 1, July 31, 2014




Problem 1.   Determine all pairs (a, b) of real numbers for whi h there exists a unique
symmetri 2 × 2 matrix M with real entries satisfying trace(M) = a and det(M) = b.
                              (Proposed by Stephan Wagner, Stellenbos h University)
Solution 1.   Let the matrix be              
                                         x z
                                      M=       .
                                         z y
The two onditions give us x + y = a and xy − z 2 = b. Sin e this is symmetri in x and
y , the matrix an only be unique if x = y . Hen e 2x = a and x2 − z 2 = b. Moreover,
if (x, y, z) solves the system of equations, so does (x, y, −z). So M an only be unique if
z = 0. This means that 2x = a and x2 = b, so a2 = 4b.
     If this is the ase, then M is indeed unique: if x + y = a and xy − z 2 = b, then

                  (x − y)2 + 4z 2 = (x + y)2 + 4z 2 − 4xy = a2 − 4b = 0,

so we must have x = y and z = 0, meaning that
                                                      
                                       a/2 0
                                    M=
                                        0 a/2

is the only solution.
Solution 2. Note that trace(M) = a and det(M) = b if and only if the two eigenvalues
λ1 and λ2 of M are solutions of x2 − ax + b = 0. If λ1 6= λ2 , then
                                                                
                            λ1 0                           λ 0
                      M1 =                    and     M2 = 2
                             0 λ2                          0 λ1

are two distin t solutions, ontradi ting uniqueness. Thus λ1 = λ2 = λ = a/2, whi h
implies a2 = 4b on e again. In this ase, we use the fa t that M has to be diagonalisable
as it is assumed to be symmetri . Thus there exists a matrix T su h that
                                                     
                                         −1      λ 0
                                  M =T        ·         · T,
                                                  0 λ

however this redu es to M = λ(T −1 · I · T ) = λI , whi h shows again that M is unique.
Problem 2.     Consider the following sequen e
                    (an )∞
                         n=1 = (1, 1, 2, 1, 2, 3, 1, 2, 3, 4, 1, 2, 3, 4, 5, 1, . . . ).
                                                                               n
                                                                               P
                                                                                       ak
Find all pairs (α, β) of positive real numbers su h that          = β.  lim k=1
                                                                       n→∞ nα
                                   (Proposed by Tomas Barta, Charles University, Prague)
         Let Nn = n+1     (then aNn is the rst appearan e of number n in the sequen e)
                             
Solution.
                       2
and onsider limit of the subsequen e
         PNn           Pn                          Pn    k+1
                                                                       n+2
                                                                                 1 3
            k=1 ak       k=1 1 + · · · + k           k=1                            n (1 + 2/n)(1 + 1/n)
bNn :=             =         n+1 α
                                           =             2
                                                      n+1 α
                                                                 =     n+1
                                                                          3
                                                                              α = 6
                                                                                                         .
            Nnα                                                                    (1/2)α n2α (1 + 1/n)α
                                                                           
                                 2                     2                 2

We an see that lim bNn is positive and nite if and only if α = 3/2. In this ase the
                  n→∞√                                         √
limit is equal to β = 32 . So, this pair (α, β) = ( 23 , 32 ) is the only andidate for solution.
We will show onvergen e of the original sequen e for these values of α and β .
   Let N be a positive integer in [Nn +1, Nn+1 ], i.e., N = Nn +m for some 1 ≤ m ≤ n+1.
Then we have                                n+2
                                                     m+1
                                                           
                                                      3
                                                         +    2
                                            bN =              3/2
                                                     n+1
                                                         
                                                       2
                                                           + m
whi h an be estimated by
                                        n+2                  n+2
                                                                  + n+1
                                                                             
                                         3                    3        2
                                  n+1
                                          3/2 ≤ bN ≤          n+1 3/2
                                                                                  .
                                   2
                                         +n                      2
                                        √
Sin e both bounds onverge to 32 , the sequen e bN has the same limit and we are done.


Problem 3.             Let n be a positive integer. Show that there are positive real numbers
a0 , a1 , . . . , an su h that for ea h hoi e of signs the polynomial

                                 ±an xn ± an−1 xn−1 ± · · · ± a1 x ± a0

has n distin t real roots.
                                                   (Proposed by Stephan Neupert, TUM, Mün hen)
Solution. We pro eed by indu tion on n. The statement is trivial for n = 1. Thus
assume that we have some an , . . . , a0 whi h satisfy the onditions for some n. Consider
now the polynomials
                         P̃ (x) = ±an xn+1 ± an−1 xn ± . . . ± a1 x2 ± a0 x

By indu tion hypothesis and a0 6= 0, ea h of these polynomials has n + 1 distin t zeros,
in luding the n nonzero roots of ±an xn ± an−1 xn−1 ± . . . ± a1 x ± a0 and 0. In parti ular
none of the polynomials has a root whi h is a lo al extremum. Hen e we an hoose some
ε > 0 su h that for ea h su h polynomial P̃ (x) and ea h of its lo al extrema s we have
|P̃ (s)| > ε. We laim that then ea h of the polynomials

                       P (x) = ±an xn+1 ± an−1 xn ± . . . ± a1 x2 ± a0 x ± ε
has exa tly n + 1 distin t zeros as well. As P̃ (x) has n + 1 distin t zeros, it admits a
lo al extremum at n points. Call these lo al extrema −∞ = s0 < s1 < s2 < . . . < sn <
sn+1 = ∞. Then for ea h i ∈ {0, 1, . . . , n} the values P̃ (si ) and P̃ (si+1 ) have opposite
signs (with the obvious onvention at innity). By hoi e of ε the same holds true for
P (si ) and P (si+1 ). Hen e there is at least one real zero of P (x) in ea h interval (si , si+1 ),
i.e. P (x) has at least (and therefore exa tly) n+ 1 zeros. This shows that we have found a
set of positive reals a′n+1 = an , a′n = an−1 , . . . , a′1 = a0 , a′0 = ε with the desired properties.


Problem 4.       Let n > 6 be a perfe t number, and let n = pe11 · · · pekk be its prime
fa torisation with 1 < p1 < . . . < pk . Prove that e1 is an even number.
    A number n is perfe t if s(n) = 2n, where s(n) is the sum of the divisors of n.
                          (Proposed by Javier Rodrigo, Universidad Ponti ia Comillas)
Solution.        Suppose that e1Qis odd, ontrary to the statement.
      We know that s(n) = ki=1 (1 + pi + p2i + · · · + pei i ) = 2n = 2pe11 · · · pekk . Sin e e1 is
an odd number, p1 + 1 divides the rst fa tor 1 + p1 + p21 + · · · + pe11 , so p1 + 1 divides
2n. Due to p1 + 1 > 2, at least one of the primes p1 , . . . , pk divides p1 + 1. The primes
p3 , . . . , pk are greater than p1 + 1 and p1 annot divide p1 + 1, so p2 must divide p1 + 1.
Sin e p1 + 1 < 2p2 , this possible only if p2 = p1 + 1, therefore p1 = 2 and p2 = 3. Hen e,
6|n.
      Now n, n2 , n3 , n6 and 1 are distin t divisors of n, so
                                        n n n
                          s(n) ≥ n +     + + + 1 = 2n + 1 > 2n,
                                        2 3 6
 ontradi tion.


Remark. It is well-known that all even perfe t numbers have the form n = 2
                                                                                     p−1 (2p − 1) su h
            p
that p and 2 − 1 are primes.
                                                                           p
                                 So if e1 is odd then k = 2, p1 = 2, p2 = 2 − 1, e1 = p − 1 and

e2 = 1. If n > 6 then p > 2 so p is odd and e1 = p − 1 should be even.


Problem 5.      Let A1 A2 . . . A3n be a losed broken line onsisting of 3n line segments
in the Eu lidean plane. Suppose that no three of its verti es are ollinear, and for
ea h index i = 1, 2, . . . , 3n, the triangle Ai Ai+1 Ai+2 has ounter lo kwise orientation and
∠Ai Ai+1 Ai+2 = 60◦ , using the notation A3n+1 = A1 and A3n+2 = A2 . Prove that the
                                                               3
number of self-interse tions of the broken line is at most n2 − 2n + 1.
                                                                     2
                                                           A3
                                                 A6




                                            A4        A5
                                       A1                       A2

                                                                         (Proposed by Martin Langer)
Solution.   Pla e the broken line inside an equilateral triangle T su h that their sides are
parallel to the segments of the broken line. For every i = 1, 2, . . . , 3n, denote by xi the
distan e between the segment Ai Ai+1 and that side of T whi h is parallel to Ai Ai+1 . We
will use indi es modulo 3n everywhere.
    It is easy to see that if i ≡ j (mod 3) then the polylines Ai Ai+1 Ai+2 and Aj Aj+1 Aj+2
interse t at most on e, and this is possible only if either xi < xi+1 and xj > xj+1 or
xi < xi+1 and xj > xj+1 . Moreover, su h ases over all self-interse tions. So, the number
of self-interse tions annot ex eed number of pairs (i, j) with the property

(∗) i ≡ j   (mod 3),        and (xi < xi+1 and xj > xj+1 ) or (xi > xi+1 and xj < xj+1 ).




                            Ai+2                                             Aj+2

                     xi+2              xi+1                             Ai+2         xi+1
                                                                                       xj+1
                Ai+3                     Ai+4                      Aj
                                                                                       Aj+1
                Ai                     Ai+1
                       xi                                        Ai             xj    Ai+1
                              xi+3                                      xi

     Grouping the indi es 1, 2, . . . , 3n, by remainders modulo 3, we have n indi es in ea h
residue lass. Altogether there are 3 n2 index pairs (i, j) with i ≡ j (mod 3). We will
show that for every integer k with 1 ≤ k < n2 , there is some index i su h that the pair
(i, i + 6k) does not satisfy (∗). This is already n−1   2
                                                            pair; this will prove that there are
at most                                        
                                        n   n−1  3
                                   3      −     ≥ n2 − 2n + 1
                                        2    2   2
self-interse tions.
     Without loss of generality we may assume that x3n = x0 is the smallest among
x1 , . . . , x3n . Suppose that all of the pairs

         (−6k, 0),     (−6k + 1, 1),            (−6k + 2, 2),   ...,    (−1, 6k − 1),         (0, 6k)   (∗∗)

satisfy (∗). Sin e x0 is minimal, we have x−6k > x0 . The pair (−6k, 0) satises (∗), so
x−6k+1 < x1 . Then we an see that x−6k+2 > x2 , and so on; nally we get x0 > x6k .
But this ontradi ts the minimality of x0 . Therefore, there is a pair in (**) that does not
satisfy (∗).

                            n   n−1   3 2        
Remark. The bound 3
                            2 −    2   = 2 n − 2n + 1 is sharp.
